% GENERATED FILE. Edit canonical lessons, then run npm run generate:books.
% canonical-source-hash: fnv1a64:ac1b314c542d516b
% canonical-lessons: RU-W260-four-lines

\chapter{Writing --- Four Lines And The Comma You Cannot Hear}
\label{ch:ru-four-lines}
\hlchaptermodality{260}

\begin{chapteropening}
\textbf{By the end of this chapter:} \emph{I can write four connected lines to a named reader, joining with \ru{и}, \ru{но} or \ru{а} and giving a reason with \ru{потому} \ru{что}, with a comma in front of every joiner that takes one.}

Complete RU-W260-four-lines: four lines to one reader using \ru{и}, \ru{но} or \ru{а}, and \ru{потому} \ru{что}, then check the comma in front of each joiner.
\end{chapteropening}

\section[\texorpdfstring{\ru{потому} \ru{что} (potomú shto)}{потому что (potomú shto)}]{Four lines, and the comma you cannot hear}

\label{lesson:RU-W260-four-lines}



\begin{quote}
The last lesson asked for writing \textbf{finished on time}. This one asks for writing that \textbf{holds together}.

Pick one person who will read it.
\end{quote}

\subsection*{What to know first}
Connected writing is not \textbf{longer} writing. Four short lines can be connected; one long line can fail to be. What makes it connected is that a sentence stops standing on its own.

You hold four joiners. \textbf{\ru{и}} adds. \textbf{\ru{но}} denies. \textbf{\ru{а}} stands one thing beside another. \textbf{\ru{потому} \ru{что}} gives the reason.

And here is the Russian thing to notice. Said aloud, a join is just a word. Written down, three of these four carry a \textbf{mark}: a comma before \textbf{\ru{но}}, before \textbf{\ru{а}}, and before \textbf{\ru{потому} \ru{что}} — before \textbf{\ru{потому}}; at this level, never inside the pair.

\textbf{\ru{и}} between two halves about the same person takes none.

So the ear never hears the comma, and the examiner always sees it. It is the one piece of this lesson a reader can mark wrong without any argument.

\subsection*{Writing — connected composition}
Write \textbf{four lines} by hand to the person you picked.

At least one must use \textbf{\ru{и}}, joining two things the same person does. At least one must use \textbf{\ru{но}} or \textbf{\ru{а}}. At least one must use \textbf{\ru{потому} \ru{что}} — both words, spelled out.

Use only words you already write. This lesson adds no vocabulary; the whole exercise is the joining.

\begin{tcolorbox}[breakable,colback=teal!4,colframe=teal!35!black,title={Before you move on}]
Cover your first line and read the second alone. Does it still make sense without the first? If it does, you wrote a \textbf{list} — which is where you were one lesson ago.

Then find every joiner and look just in front of it. A comma before \textbf{\ru{но}}, \textbf{\ru{а}} and \textbf{\ru{потому}}; none before an \textbf{\ru{и}} that joins two things one person does. Each comma's tail below the line.
\end{tcolorbox}
